% !TEX root = ../main.tex
%
% Methods subsection -- Evaluation Protocol

\subsection{Evaluation Protocol}
\label{subsec:evaluation-protocol}

\textbf{Internal evaluation.} Within each training cohort, patients are split
case-wise into training, validation, and internal test sets in a 6:2:2 ratio,
stratified jointly by cohort and event status so that the death/censoring ratio
is preserved across splits. No patient, and no slide belonging to that patient,
appears in more than one split.

\textbf{External (cross-institution) evaluation.} Each model is trained
exclusively on one cohort and evaluated, without any further fine-tuning, on the
\emph{entire} other cohort, which is never seen during training or model
selection. This is a stricter test of generalization than pooling both cohorts
before splitting: pooling exposes both institutions' scanning/preparation
characteristics to the model during training, whereas external evaluation
requires the model to generalize to an institution it has never encountered.
Both training directions are evaluated (TCGA-PAAD $\rightarrow$ CPTAC-PDAC and
CPTAC-PDAC $\rightarrow$ TCGA-PAAD).

\textbf{Metrics.} Four survival metrics are reported:

\begin{enumerate}
  \item \emph{Harrell's C-index}: among all comparable patient pairs (an
    earlier death compared against a patient observed at least that long), the
    fraction whose predicted risk ordering agrees with the observed survival
    ordering, with tied risk scores counted as half-concordant.
  \item \emph{Hazard ratio (HR)} with 95\% confidence interval: patients are
    split into high- and low-risk groups at the median predicted risk score,
    and a Cox model with group membership as its sole covariate is fit to
    obtain the HR.
  \item \emph{Log-rank test $p$-value}: computed between the Kaplan--Meier
    curves of the two median-split risk groups.
  \item \emph{Time-dependent AUC} (Uno's IPCW-weighted cumulative/dynamic
    AUC): measures how well the risk score discriminates patients who have
    died by a given time point ($t$) from those still observed beyond it,
    evaluated at $t = $ 12, 24, and 36 months; the censoring-time distribution
    used for inverse-probability weighting is estimated from the training
    cohort, following standard convention.
\end{enumerate}

\textbf{Aggregation across runs.} Following the repeated-run protocol
(multiple seeds, both external training directions, and pooled 5-fold internal
evaluation), all four metrics are reported as mean $\pm$ SD across runs rather
than from a single training run, given that random initialization alone
produces C-index variation on the same order as differences between competing
model variants.
